\uuid{WtTm}
\exo7id{3883}
\titre{exo7 3883}
\auteur{quercia}
\organisation{exo7}
\datecreate{2010-03-11}
\isIndication{false}
\isCorrection{true}
\chapitre{Continuité, limite et étude de fonctions réelles}
\sousChapitre{Etude de fonctions}
\cpge{mpsi-mp2i.fonctions-valeurs-reelles-complexes.a, pcsi.techniques-calcul-fonctions.a, ptsi.techniques-calcul-fonctions.a, mpsi-mp2i.limites-continuite.a, pcsi.limites-continuite.a, ptsi.limites-continuite.a, mpsi-mp2i.limites-continuite.b, pcsi.limites-continuite.b, ptsi.limites-continuite.b, bcpst1.fonctions-reelles-usuelles.a, bcpst1.limites-continuite.a, bcpst1.limites-continuite.c}
\module{Analyse}
\niveau{L1}
\difficulte{}

\contenu{
\question{
Soit $f : {[0,1]} \to {[0,1]}$ croissante. Montrer qu'il existe $x \in {[0,1]}$
tel que $f(x) = x$.

(\'Etudier $A = \{ x \in [0,1] \text{ tq } f(x) \le x \}$)
}
\reponse{
$\alpha = \inf(A)  \Rightarrow  f(\alpha^+) \le \alpha \text{ et } f(\alpha^-) \ge \alpha$.
}
}
